% !TeX program = lualatex
% Compile twice: lualatex 123.tex
% All references and prose are embedded; no BibTeX database or external inputs.
\documentclass[11pt,a4paper]{article}
\usepackage[margin=24mm,headheight=14pt]{geometry}
\usepackage{fontspec}
\setmainfont{Latin Modern Roman}
\setsansfont{Latin Modern Sans}
\setmonofont{Latin Modern Mono}
\usepackage{amsmath,amssymb,mathtools}
\usepackage{unicode-math}
\setmathfont{Latin Modern Math}
\usepackage{microtype}
\usepackage{enumitem,xurl,needspace}
\usepackage[dvipsnames]{xcolor}
\usepackage{hyperref}
\hypersetup{unicode=true,colorlinks=true,linkcolor=MidnightBlue,urlcolor=MidnightBlue,
 pdftitle={500 Mathematical Problem Statements},pdfauthor={Source-annotated compilation},bookmarksnumbered=true}
\usepackage{bookmark}
\usepackage{tocloft}
\setlength{\cftsecnumwidth}{3em}
\cftsetpnumwidth{2.5em}
\cftsetrmarg{4em}
\usepackage{titlesec}
\titleformat{\section}{\Large\bfseries\raggedright}{\thesection}{0.7em}{}
\usepackage{fancyhdr}
\pagestyle{fancy}\fancyhf{}
\fancyhead[L]{\small\sffamily Mathematical problem statements}
\fancyhead[R]{\small Source edition 22}
\fancyfoot[C]{\thepage}
\setlength{\parindent}{0pt}
\setlength{\parskip}{5pt plus 1pt}
\setlength{\emergencystretch}{3em}
\setcounter{tocdepth}{1}
\setcounter{secnumdepth}{1}
\setlist[itemize]{leftmargin=3em,itemsep=5pt,topsep=4pt,parsep=0pt}
\allowdisplaybreaks
\newcommand{\N}{\mathbb{N}}
\newcommand{\Z}{\mathbb{Z}}
\newcommand{\Q}{\mathbb{Q}}
\newcommand{\R}{\mathbb{R}}
\newcommand{\C}{\mathbb{C}}
\newcommand{\F}{\mathbb{F}}
\newcommand{\A}{\mathbb{A}}
\newcommand{\PP}{\mathbb P}
\newcommand{\E}{\mathbb{E}}
\newcommand{\eps}{\varepsilon}
\newcommand{\dd}{\,\mathrm d}
\newcommand{\sref}[2]{\hyperlink{source:#1:#2}{[S#2]}}
\newcommand{\eref}[2]{\hyperlink{extra:#1:#2}{[E#2]}}
% Turn the next switch off for a shorter reading copy. The quotations preserve
% every exactTarget field, including qualifications omitted from a short summary.
\newif\ifshowcatalogue
\showcataloguetrue
\newcommand{\cataloguescope}[1]{\ifshowcatalogue
 \paragraph{Exact scope in the supplied catalogue.}
 \begin{quote}\small #1\end{quote}\fi}
\hypersetup{pdftitle={Problem 123 - Positive scalar curvature and asphericity},pdfauthor={Open Problems Math research notebook}}
\fancyhead[L]{\small\sffamily Open Problems Math \textbar\ Research notebook}
\fancyhead[R]{\small\sffamily 123 \textbar\ 27 September 2026}
\numberwithin{equation}{section}
\begin{document}
\setcounter{section}{122}
\setlength{\parskip}{3pt plus 1pt}
\setlist[itemize]{itemsep=3pt,topsep=3pt}
\titlespacing*{\paragraph}{0pt}{1.5ex plus .5ex minus .2ex}{1em}
% ================================================================
% problemId: problem.aspherical-manifolds-admit-no-metric-of-positive-scalar-curvature
% source releaseRank: 123; source releaseStatus: open_with_solved_subcases
% exactTarget SHA-256: 9ef5b78df2371c1583410b14cb5b51ab2d4fad676eebd22da7f5d039ec50c8f4
\clearpage
\section{\texorpdfstring{Aspherical Manifolds Admit No Metric of Positive Scalar Curvature}{Aspherical Manifolds Admit No Metric of Positive Scalar Curvature}}\label{123}
\subsection{Definitions and mathematical statement}
A closed manifold is compact without boundary, and asphericity means its universal cover is contractible. The scalar curvature of a Riemannian metric is the trace of its Ricci curvature. The target is
\[
 M\text{ closed and aspherical}\quad\Longrightarrow\quad
 \nexists g\text{ on }M\text{ with }\operatorname{Scal}_g(x)>0\text{ for every }x.
\]
Strict pointwise positivity is intended. A prohibition only on positive \emph{sectional} curvature is weaker, and additional spin or fundamental-group assumptions would restrict the conjecture.
\par\smallskip\textit{Formulation and concept references:} \sref{123}{1}.
\cataloguescope{Prove or disprove that no closed aspherical manifold admits a Riemannian metric of positive scalar curvature.}
\subsection{Short English statement}
Does asphericity rule out every everywhere-positive scalar-curvature metric on a closed manifold?
% BEGIN RESEARCH ADDITION 123
\subsection{Research attempt}

\paragraph{Editorial source note.}
The original source entry is retained unchanged.  Its URL identifies a
paper by \emph{Shuli Chen}, Jianchun Chu, and Jintian Zhu, not by Simone
Cecchini, Jianchun Chu, and Jintian Zhu.  The corrected bibliographic
entry is added below. \eref{123}{1}

\paragraph{Current frontier.}
The December 2025 version of that paper records the closed aspherical
obstruction through dimension five, and proves stronger connected-sum
and noncompact-domination results in dimensions three, four, and five.
It explicitly distinguishes these from the general higher-dimensional
problem.  Its introduction also explains why asphericity alone need not
supply the homology classes used in elementary dimension descent.
\eref{123}{1}, introduction and Theorems 1.2 and 1.8.
The spinorial curvature obstruction used below is standard; Listing's
formulation makes the operator-norm dependence of the twisted
Bochner formula explicit. \eref{123}{2}

\paragraph{Attack routes.}
The index route needs a nonzero index that survives twisting by bundles
of arbitrarily small curvature.  Contractibility of the universal cover
does not directly construct those bundles.  The minimal-hypersurface
route needs essential smooth hypersurfaces and an inductively preserved
topological obstruction.  A counterexample route would construct a
positive-scalar-curvature metric without accidentally losing asphericity
through surgery.  We pursue the first two estimates and then identify
the topology that neither estimate provides.  Spin is an additional
hypothesis of one route, not of the original problem.

\paragraph{Attempt I: the exact analytic input for a spin obstruction.}
Let $M^n$, $n\geq2$, be closed, spin, and even-dimensional, with a hypothetical
metric $g$ of scalar curvature at least $\kappa>0$.  The positive
constant exists because $M$ is compact and the curvature is strictly
positive.  For a Hermitian bundle $E$ with unitary connection, the
twisted spin Dirac operator satisfies
\begin{equation}\label{123:dirac}
 D_E^2=\nabla^*\nabla+\frac{R_g}{4}+\mathcal R_E,
 \qquad
 \mathcal R_E=\sum_{i<j}c(e_i)c(e_j)F_E(e_i,e_j).
\end{equation}
Here $c$ is Clifford multiplication and $F_E$ is curvature.  This is
the twisted Bochner--Lichnerowicz identity used in \eref{123}{2},
Sections 2 and 3.3.  Define $\|F_E\|$ by the supremum of the operator
norm on orthonormal pairs of tangent vectors.  Since unit Clifford
multiplications have norm one,
\[
 \|\mathcal R_E\|\leq\binom n2\|F_E\|.
\]
In particular this bound does not acquire an extra factor equal to
$\operatorname{rank}E$.

If $\|F_E\|<\kappa/(8\binom n2)$, integration of
\eqref{123:dirac} for a smooth spinor $s$ gives
\[
 \|D_Es\|_2^2\geq\|\nabla s\|_2^2+\frac\kappa8\|s\|_2^2.
\]
Both chiral kernels vanish, hence $\operatorname{index}D_E^+=0$.
Therefore the following extra geometric condition rules out $g$:
\begin{equation}\label{123:bundle}
 \text{for every $\eta>0$ there is $E$ with }
 \|F_E\|<\eta\text{ and }\operatorname{index}D_E^+\ne0.
\end{equation}
This implication is completely proved.  The issue is whether the
appropriate topological construction supplies \eqref{123:bundle} for
the manifold under consideration; the implication does not prove its
hypothesis.  Nontriviality of $E$ alone is not enough, since it need not
pair nontrivially with the Dirac index.

For an odd-dimensional spin manifold, the same conditional obstruction
can be applied to $M\times S^1$: a product with a flat circle preserves
strictly positive scalar curvature and gives even dimension.  It still
requires the corresponding nonzero twisted indices on the product.
For a nonspin manifold, passing to the contractible universal cover
makes that cover spin but noncompact.  The closed Fredholm-index
argument just given then does not apply.  Existence of a finite spin
cover has not been assumed or proved.  An orientable double cover does
not in general remove the spin obstruction.

\paragraph{Attempt II: what stable minimal hypersurfaces do provide.}
Suppose $\Sigma^m\subset M^{m+1}$ is a closed smooth two-sided stable
minimal hypersurface, and $R_M\geq\kappa>0$.  Stability and the Gauss
identity give, for every smooth real $\phi$ on $\Sigma$,
\[
 \int_\Sigma|\nabla\phi|^2\geq
 \int_\Sigma(\operatorname{Ric}_M(\nu,\nu)+|II|^2)\phi^2,
 \quad
 R_\Sigma=R_M-2\operatorname{Ric}_M(\nu,\nu)-|II|^2.
\]
Eliminating the normal Ricci term yields the useful lower bound
\begin{equation}\label{123:stable}
 \int_\Sigma\bigl(2|\nabla\phi|^2+R_\Sigma\phi^2\bigr)
 \geq\int_\Sigma(R_M+|II|^2)\phi^2
 \geq\kappa\int_\Sigma\phi^2.
\end{equation}
This reconstructs the scalar-curvature part of the dimension-descent
mechanism underlying the source's geometric discussion.
\eref{123}{1}, Sections 1.1 and 2.6.

For $m\geq3$, the conformal Laplacian has gradient coefficient
$a_m=4(m-1)/(m-2)>2$.  Its first eigenvalue is positive by
\eqref{123:stable}, and its first eigenfunction $u$ can be chosen
strictly positive.  The scalar curvature of the conformally changed
metric $u^{4/(m-2)}g_\Sigma$ is
\[
 u^{-(m+2)/(m-2)}(-a_m\Delta_\Sigma u+R_\Sigma u)>0.
\]
For an orientable surface $m=2$, using $\phi=1$ in
\eqref{123:stable} gives $\int R_\Sigma>0$, hence positive Euler
characteristic by Gauss--Bonnet.  A connected such surface is a sphere.
Thus a smooth stable hypersurface inherits a precise positive-scalar-
curvature conclusion, with no unproved sign in the second variation.

The desired induction would now need this $\Sigma$ to carry an
obstruction contradicting its positive curvature, for example to be
aspherical.  That conclusion is absent.  An embedded hypersurface of
an aspherical manifold need not be aspherical: the boundary of a small
coordinate ball is a sphere.  This example does not claim that the
sphere is stable minimal in a hypothetical positive-scalar-curvature
metric; it refutes the purely topological inheritance assertion that
an attempted induction would otherwise use.  Essentiality, existence,
and stability must all be supplied together.

Even constructing essential hypersurfaces by dualizing first cohomology
cannot cover all aspherical manifolds.  The formulation source notes
aspherical homology-sphere examples in dimension four.  A strategy
requiring nonzero first Betti number has already restricted the target.
In higher dimensions a minimizing hypersurface may also have a singular
set; the smooth conformal argument cannot simply be iterated through
that set without a separate theorem. \eref{123}{1}, introduction.

\paragraph{A counterexample-oriented check.}
Attaching handles or taking connected sums is a familiar way to alter
a metric while retaining positive scalar curvature under suitable
surgery hypotheses.  It is not a way to preserve asphericity for free.
For example, the connected sum with $S^2\times S^{n-2}$, in dimensions
where this factor is simply connected, introduces an explicit
higher-homotopy obstruction detectable in the universal cover.
No surgery construction is presented here with a verified contractible
universal cover and an everywhere-positive scalar-curvature metric.
Thus these metric modifications do not furnish a counterexample to the
catalogue's exact class.

\paragraph{Stress test.}
The index estimate uses curvature in operator norm, not a trace bound
that could deteriorate as the twisting rank grows.  The inequality is
applied on a closed manifold, where the ordinary chiral index is defined;
it is not extended to a noncompact cover without Fredholm hypotheses.
The minimal-hypersurface estimate assumes smoothness, two-sidedness,
and stability explicitly, and establishes scalar, not sectional,
curvature.  Neither spin nor an abundance of degree-one cohomology has
been added to the universal statement.

The companion symbolic checks verify the elimination of the Ricci term
and the inequalities between the conformal-Laplacian coefficients.
They test algebraic signs only; they do not construct missing bundles,
essential hypersurfaces, or singularity-removal arguments.

\paragraph{Outcome.}
\textbf{Outcome: Exploratory approach with unresolved gap.}

Both analytic mechanisms have been pushed to precise sufficient
hypotheses.  Small-curvature nonzero-index twists give a contradiction
in the indicated spin setting, and smooth stable hypersurfaces inherit
positive scalar curvature in the stated sense.  The general aspherical
topology has not been shown to supply the objects required by either
mechanism, and the nonspin and high-dimensional difficulties remain.
No extension of the established dimension range or new counterexample
is claimed.
% END RESEARCH ADDITION 123
\subsection{Sources}
\begin{itemize}\raggedright
\item[\textnormal{[S1]}]\hypertarget{source:123:1}{} Simone Cecchini, Jianchun Chu, and Jintian Zhu, Positive scalar curvature metrics and aspherical summands, arXiv:2312.04698v3 (2025), introduction.
\url{https://arxiv.org/html/2312.04698}.
{\small\textit{Catalogue formulation source.}}
% BEGIN RESEARCH REFERENCES 123
\item[\textnormal{[E1]}]\hypertarget{extra:123:1}{} Shuli Chen, Jianchun Chu, and Jintian Zhu.
\emph{Positive scalar curvature metrics and aspherical summands}, arXiv:2312.04698v3, 18 December 2025. Introduction, Theorems 1.2 and 1.8, and the geometric preliminaries. Corrects the first author in the retained catalogue reference.
\url{https://arxiv.org/html/2312.04698v3}.
{\small\textit{Consulted 27 September 2026 for the indicated result or scope.}}
\item[\textnormal{[E2]}]\hypertarget{extra:123:2}{} Mario Listing.
\emph{Scalar curvature and vector bundles}, arXiv:1202.4325v1, 20 February 2012. Sections 2 and 3.3, including the twisted Dirac curvature term and its operator-norm estimate.
\url{https://arxiv.org/html/1202.4325v1}.
{\small\textit{Consulted 27 September 2026 for the indicated result or scope.}}
% END RESEARCH REFERENCES 123
\end{itemize}


\end{document}
